\documentclass[a4paper,11pt]{article}
\usepackage[T2A]{fontenc}
\usepackage[utf8]{inputenc}
\usepackage[russian]{babel}
\usepackage{amsmath,amssymb,mathtools}
\usepackage{cmbright}
\renewcommand{\familydefault}{\sfdefault}
\usepackage{icomma}
\usepackage[a4paper,left=18mm,right=18mm,top=11mm,bottom=11mm]{geometry}
\usepackage{microtype}
\usepackage{booktabs,tabularx}
\usepackage{xcolor}
\usepackage{hyperref}
\definecolor{accent}{HTML}{15747A}
\definecolor{darkaccent}{HTML}{164C55}
\definecolor{textgray}{HTML}{283238}
\definecolor{muted}{HTML}{65747A}
\definecolor{rulegray}{HTML}{CADADC}
\color{textgray}
\hypersetup{hidelinks,pdftitle={Главные формулы тригонометрии. Ярослав Гаврилов},pdfauthor={Лёвин Артём Александрович}}
\setlength{\parindent}{0pt}
\setlength{\parskip}{0pt}
\setlength{\arraycolsep}{3pt}
\setlength{\abovedisplayskip}{0pt}\setlength{\abovedisplayshortskip}{0pt}
\setlength{\belowdisplayskip}{0pt}\setlength{\belowdisplayshortskip}{0pt}
\setlength{\jot}{1.5pt}
\newcommand{\R}{\mathbb{R}}
\newcommand{\Z}{\mathbb{Z}}
\newcommand{\sect}[1]{\par\vspace{0.8mm}{\normalsize\bfseries\color{darkaccent}#1}\par\vspace{0.25mm}{\color{rulegray}\rule{\linewidth}{0.4pt}}\par\vspace{0.1mm}}
\newcommand{\mathline}[1]{\[\displaystyle #1\]}
\begin{document}
\thispagestyle{empty}
\begin{center}
{\small\bfseries\color{accent}ПРОФИЛЬНАЯ МАТЕМАТИКА \quad\textbullet\quad ПАМЯТКА}\par
\vspace{0.9mm}
{\fontsize{22}{26}\selectfont\bfseries\color{darkaccent}Главные формулы тригонометрии}\par
\vspace{0.4mm}
{\small\color{muted}Ярослав Гаврилов \quad\textbullet\quad 12.10.2026}
\end{center}
\vspace{0.75mm}
\noindent{\color{accent}\rule{\linewidth}{0.9pt}}
\vspace{0.7mm}
{\normalsize\bfseries\color{darkaccent}Главные принципы решения}\par\vspace{1.2mm}
\begin{minipage}[t]{0.47\linewidth}
\footnotesize
\textbf{Тригонометрические уравнения:} стремитесь к одному углу и одной функции.\par\vspace{0.8mm}
\textbf{Разные аргументы:} $x$ и $2x$ --- проверьте формулы двойного угла.\par\vspace{0.8mm}
\textbf{Произведение равно нулю:} приравняйте к нулю каждый множитель.\par
\end{minipage}\hfill
\begin{minipage}[t]{0.47\linewidth}
\footnotesize
\textbf{Деление:} не делите на $\sin x$ или $\cos x$ без проверки нулевого случая.\par\vspace{0.8mm}
\textbf{Знаки и ограничения:} проверяйте ОДЗ, четверть и знак после извлечения корня.\par\vspace{0.8mm}
\textbf{Ответ:} сначала все серии корней ($n\in\Z$), затем отбор на промежутке.\par\vspace{0.8mm}
\textbf{Неравенства:} выбирайте подходящие дуги на единичной окружности.\par
\end{minipage}
\vspace{2mm}


\sect{1. Основное тождество, тангенс и котангенс}
\begin{align*}
\sin^2\alpha+\cos^2\alpha&=1, \\
\tg\alpha&=\frac{\sin\alpha}{\cos\alpha}\quad(\cos\alpha\ne0),\\
\ctg\alpha&=\frac{\cos\alpha}{\sin\alpha}\quad(\sin\alpha\ne0).
\end{align*}

\sect{2. Сумма и разность углов}
\begin{align*}
\sin(\alpha+\beta)&=\sin\alpha\cos\beta+\cos\alpha\sin\beta,\\
\sin(\alpha-\beta)&=\sin\alpha\cos\beta-\cos\alpha\sin\beta,\\[0.4mm]
\cos(\alpha+\beta)&=\cos\alpha\cos\beta-\sin\alpha\sin\beta,\\
\cos(\alpha-\beta)&=\cos\alpha\cos\beta+\sin\alpha\sin\beta.
\end{align*}


\noindent\begin{minipage}[t]{0.49\linewidth}
\sect{3. Двойной угол}
\vspace{1mm}
\begin{align*}
\sin 2\alpha&=2\sin\alpha\cos\alpha,\\
\cos 2\alpha&=\cos^2\alpha-\sin^2\alpha\\
&=1-2\sin^2\alpha\\
&=2\cos^2\alpha-1.
\end{align*}
\end{minipage}\hfill
\begin{minipage}[t]{0.49\linewidth}
\sect{4. Табличные значения}
\vspace{0.3mm}
\centering
\renewcommand{\arraystretch}{1.04}
\begin{tabular}{@{}lccccc@{}}
\toprule
$\alpha$ & $0$ & $\dfrac{\pi}{6}$ & $\dfrac{\pi}{4}$ & $\dfrac{\pi}{3}$ & $\dfrac{\pi}{2}$\\
\midrule
$\sin\alpha$ & $0$ & $\dfrac12$ & $\dfrac{\sqrt2}{2}$ & $\dfrac{\sqrt3}{2}$ & $1$\\
$\cos\alpha$ & $1$ & $\dfrac{\sqrt3}{2}$ & $\dfrac{\sqrt2}{2}$ & $\dfrac12$ & $0$\\
$\tg\alpha$ & $0$ & $\dfrac{\sqrt3}{3}$ & $1$ & $\sqrt3$ & ---\\
\bottomrule
\end{tabular}
\end{minipage}

\sect{5. Простейшие тригонометрические уравнения}
\begin{align*}
\sin x=a\quad &\Longrightarrow\quad
\begin{cases}
x=\arcsin a+2\pi n,\\
x=\pi-\arcsin a+2\pi n,
\end{cases}\\
\cos x=a\quad &\Longrightarrow\quad x=\pm\arccos a+2\pi n,\\
\tg x=a\quad &\Longrightarrow\quad x=\arctg a+\pi n.
\end{align*}
{\small\color{muted}Здесь $n\in\Z$; для синуса и косинуса $|a|\le1$. Для $\tg x$ допустимо любое $a\in\R$.}

\vfill
{\color{accent}\rule{\linewidth}{0.8pt}}\par\vspace{2mm}
{\small\color{muted}Лёвин Артём Александрович эксклюзивно для Ярослава Гаврилова}
\end{document}
